\begin{table}[htbp]
\centering
\begin{threeparttable}
\caption{CHFS 2011: Information Channels and Inflation Expectations}
\label{tab:chfs_mechanism_info_channel}
{\small
\begin{tabular}{lccccc}
\toprule
 & (1) & (2) & (3) & (4) & (5) \\
 & Infl.\ exp. & Extremity & Over-predict & Risk placebo & Horse race \\
\midrule
\multicolumn{6}{l}{\textit{Panel A: Information Channel}} \\
\addlinespace
High media (TV/internet/newspaper) & $-0.005$ & $-0.001$ & $-0.004$ & & $-0.002$ \\
 & $(0.022)$ & $(0.015)$ & $(0.012)$ & & $(0.022)$ \\
 & $[0.820]$ & $[0.947]$ & $[0.739]$ & & $[0.928]$ \\
\addlinespace
\multicolumn{6}{l}{\textit{Panel B: Risk Preference Placebo}} \\
\addlinespace
Risk attitude & & & & $-0.010$ & $-0.010$ \\
 & & & & $(0.008)$ & $(0.008)$ \\
 & & & & $[0.211]$ & $[0.211]$ \\
\addlinespace
Observations & 8,281 & 8,281 & 8,281 & 8,173 & 8,173 \\
Province FE & Yes & Yes & Yes & Yes & Yes \\
Controls & Yes & Yes & Yes & Yes & Yes \\
\bottomrule
\end{tabular}
}
\begin{tablenotes}[flushleft]
\footnotesize
\item \textit{Notes:} Columns (1), (4), and (5) use inflation expectation (1--5 scale, higher = more inflationary) as the dependent variable. Column (2) uses absolute deviation from province mean expectation. Column (3) uses an indicator for expecting prices to rise a lot. Controls include log income and an urban indicator. Heteroskedasticity-robust standard errors are in parentheses; p-values are in brackets. $^{*}p<0.10$; $^{**}p<0.05$; $^{***}p<0.01$.
\end{tablenotes}
\end{threeparttable}
\end{table}
